\section{Kết luận, hạn chế và hướng phát triển}

\subsection{Kết luận}

Báo cáo cho thấy bài toán forecasting trên FreshRetailNet-50K nên được framing như bài toán stockout-aware demand forecasting thay vì raw sales forecasting. Khi không xử lý stockout, observed-sales LightGBM đạt WAPE 24.42\%. Sau khi recover latent demand và train trên recovered demand, WAPE giảm xuống 20.28\%. Seasonal naive 7-day là baseline mạnh với WAPE 19.85\%. Mô hình cuối hybrid seasonal-ML đạt WAPE 19.67\%.

Kết quả này có hai thông điệp chính. Thứ nhất, xử lý stockout là cần thiết vì observed sales có thể understates demand. Thứ hai, seasonality tuần là pattern nền cực mạnh; một mô hình forecasting tốt nên kết hợp seasonality này thay vì bỏ qua.

\subsection{Hạn chế}

Recovered demand vẫn là proxy được ước lượng, không phải ground truth demand thật. Imputation ở cấp hourly row-level nhiễu vì dữ liệu thưa và nhiều zero. Test window chỉ kéo dài 14 ngày nên diagnostics non-overlap còn ít điểm. Ngoài ra, báo cáo chưa trực tiếp tối ưu inventory policy như safety stock, reorder point, service level hoặc lost-sales cost.

\subsection{Hướng phát triển}

Nếu có thêm thời gian, hướng phát triển đầu tiên là kiểm tra pipeline trên nhiều sample hoặc full data. Hướng thứ hai là thử các mô hình imputation chuyên biệt như SAITS, TimesNet hoặc DLinear ở quy mô nhỏ để so sánh với LightGBM recovery. Hướng thứ ba là chuyển từ forecast accuracy sang inventory impact: service level, fill rate, lost sales và holding cost. Đây là bước quan trọng để biến forecast thành quyết định vận hành.
